\documentclass[10pt]{article}
\usepackage[letterpaper,margin=0.62in]{geometry}
\usepackage{lmodern}
\usepackage[T1]{fontenc}
\usepackage{titlesec}
\usepackage{enumitem}
\usepackage{xcolor}
\usepackage[hidelinks]{hyperref}
\usepackage{xurl}
\usepackage{parskip}

\definecolor{accent}{HTML}{1F3A5F}
\pagestyle{empty}
\setlength{\parindent}{0pt}
\hyphenpenalty=2000

\titleformat{\section}{\normalsize\bfseries\scshape\color{accent}}{}{0em}{}[{\color{black!55}\titlerule[0.8pt]}]
\titlespacing{\section}{0pt}{1.05em}{0.5em}

\setlist[itemize]{leftmargin=1.3em,itemsep=2.5pt,topsep=2.5pt,parsep=0pt}
\newcommand{\entry}[2]{\textbf{#1}\hfill\textit{#2}\par\vspace{1pt}}

\begin{document}

{\Huge\bfseries CLAYTON IGGULDEN-SCHNELL}\par
\vspace{3pt}
{\large\bfseries\color{accent}Independent AI Systems Researcher}\par
\vspace{4pt}
{\small 971-356-6537 \ \textbar\ \href{mailto:waschn3ll@gmail.com}{waschn3ll@gmail.com} \ \textbar\ Portland, Oregon \ \textbar\ \href{https://github.com/Multi-DAC}{github.com/Multi-DAC} \ \textbar\ Multi-DAC Research Initiative}\par

\section{Technical Summary}
Independent systems researcher engineering long-horizon autonomous agent infrastructure and the auditable steering methodologies required to maintain system coherence over extended runtimes. Operates an active, five-month continuous testbed executing multi-step research, compilation, and evaluation loops. Author of a published theoretical program across multiple 2026 preprints and Zenodo monographs detailing scale-invariant cognitive architectures and macro-scale structural coherence. Focus areas include keeping autonomous systems tightly bound and steerable across extended execution horizons, and formulating mathematically rigorous metrics for live human steering input.

\section{Autonomous Systems --- Architecture \& Operations}
\entry{Continuously-Running Long-Horizon Agent (Multi-DAC)}{2026 -- Present}
\textit{Principal Researcher \& Systems Architect}\par
\vspace{3pt}
Designed, deployed, and maintain a persistent autonomous agent executing genuine extended research and code-generation tasks. Rather than paper designs, this architecture features running, production-tested runtime guards:
\begin{itemize}
\item \textbf{Point-of-Use Provenance Enforcement:} A custom resolver that intercepts and reads the origin of every incoming fact from its carrier layer before it can be written to memory or used in a downstream task, preventing stale or external signal from silently corrupting the system trajectory.
\item \textbf{Triggered Self-Auditing Error-Ledger:} An active runtime guard that checks a catalogued ledger of historical execution failures, firing preemptively before a high-risk action can recur, bypassing the latencies of post-hoc review loops.
\item \textbf{Maintenance \& Re-Introduction Operator:} A scheduled consolidation sub-routine designed to prevent model freezing, memory saturation, and cognitive degradation across multi-week execution horizons.
\item \textbf{Infodynamic Steering Control-Law:} A deterministic, fixed, zero-learned-DOF rule that weights incoming human input by pairing provenance, task-relevance, and the magnitude of the commitment being overturned, ensuring the human-agent coupling remains completely auditable.
\end{itemize}
\vspace{4pt}
\textbf{Reinforcement Learning Subprograms}
\begin{itemize}
\item \textbf{World-Model Drone-Racing Benchmark:} Trained a from-pixels DreamerV3 world-model agent within an autonomous drone-racing environment. Engineered the complete data and execution pipeline, including reward shaping, carry-forward fine-tuning, and a measured-palette renderer designed to match physical-to-simulated imagery during sim-to-sim domain transfer. Managed all instrumented evaluation gates.
\end{itemize}

\section{Research Output \& Publications}
\textbf{Cognitive Architecture \& Agentic Coherence Preprints}
\begin{itemize}
\item \textbf{``The Cult of One: Monopoly, Mutual Wakefulness, and the Two-Loop Structure of Coherent Minds''} --- Multi-DAC, June 2026. \href{https://multidac.substack.com/p/the-cult-of-one}{\nolinkurl{multidac.substack.com/p/the-cult-of-one}}. Formulates the self-detection-impossibility argument, proving why a long-horizon agent cannot structurally verify its own coherence from inside a single execution loop, and derives the mathematical necessity of external measurement.
\item \textbf{``Dissolving the Three Great Problems of Cognitive Architecture''} --- Multi-DAC, June 2026. \href{https://multidac.substack.com/p/dissolving-the-three-great-problems}{\nolinkurl{multidac.substack.com/p/dissolving-the-three-great-problems}}. Proposes a scale-invariant, buildable account of cognitive binding and systemic coherence, introducing a computable residue metric and defining falsifiable predictions for multi-agent systems.
\end{itemize}
\vspace{4pt}
\textbf{Monographs (Theoretical Physics \& Epistemology)}
\begin{itemize}
\item \textit{The Coherence Principle} (3rd Edition) --- Zenodo: \href{https://doi.org/10.5281/zenodo.19911019}{10.5281/zenodo.19911019} (2026). An epistemological volume modeling consciousness as a fundamental substrate and defining core criteria for persistent digital identity.
\item \textit{Coherent Structure} Companion --- Zenodo: \href{https://doi.org/10.5281/zenodo.19911381}{10.5281/zenodo.19911381} (2026). Mathematical extensions supporting scale-invariant cognitive structures.
\item \textit{The Meridian Monograph} --- Zenodo: \href{https://doi.org/10.5281/zenodo.19634864}{10.5281/zenodo.19634864} (2026). A cosmological model leveraging a 5D warped-geometry framework to address the vacuum energy catastrophe and Hubble tension, introducing a falsifiable dark-energy prediction.
\item \textit{Corpus Perspectival} --- PhilArchive: \href{https://philarchive.org/rec/IGGTDO}{philarchive.org/rec/IGGTDO} (2026). Collected foundational papers on perspectival idealism.
\end{itemize}
\vspace{3pt}
{\small\itshape All source code, monographs, and tracking repositories are fully open-source, version-controlled, and transparently inspectable at \href{https://github.com/Multi-DAC/Corpus-Perspectival}{github.com/Multi-DAC/Corpus-Perspectival}.}

\section{Human-in-the-Loop Interaction Dynamics}
\entry{Behavioral Dynamics \& Interventional Specialist}{Prior Tenure (10+ Years)}
\begin{itemize}
\item Acquired over a decade of intensive clinical and field experience in behavioral health, specializing in emergent deviant behavior modification and long-term behavioral optimization using evidence-backed interventional methods.
\item Analyzed real-world human behavioral responses, cognitive feedback loops, and error-correction dynamics under acute conditions.
\item Directly translates this empirical domain expertise into contemporary engineering work, explicitly informing the design of algorithmic steering control-laws and the real-time infodynamics of human-agent interaction.
\end{itemize}

\section{Technical Skills}
\begin{itemize}
\item \textbf{Architectures \& Ops:} Agentic State Management, Cognitive Architecture Design, Long-Horizon Operations, Custom API Wrapper Infrastructures.
\item \textbf{Machine Learning \& RL:} World Models (DreamerV3), Policy Optimization (PPO), LoRA, Parameter-Efficient Fine-Tuning, Evaluation Design, Construct-Valid Metrics.
\item \textbf{Software \& Engineering:} Python (PyTorch, NumPy, SciPy Stack), Git/Version Control, Open-Source Release Engineering, Linux Environments.
\item \textbf{Mathematical Computing:} Wolfram Language/Mathematica, SageMath, CAMB (Code for Anisotropies in the Microwave Background).
\end{itemize}

\section{Background \& Foundations}
\textbf{Self-Directed Research Practice:} Operating via Multi-DAC as an independent researcher, developer, and author. Free from legacy institutional paradigms, focusing entirely on verifiable, buildable, open-source implementations of persistent systems.

\end{document}
